% ---- pg-343 （书页 331） ----

\noindent $Z_4=$ 扩张血压(DIAP)

\noindent $Z_5=$ QRS 波测量值(QRS)

\noindent (a) 只用第一个响应 $Y_1$ 作回归分析.

\noindent (i) 提出并拟合适当的线性回归模型.

\noindent (ii) 分析残差.

\noindent (iii) 当 $z_1=1$, $z_2=1\,200$, $z_3=140$, $z_4=70$, 及 $z_5=85$ 时对 $Y_1$ 构造 95\% 预测区间.

\noindent (b) 用第二个响应 $Y_2$ 重复 (a).

\noindent (c) 用两个响应 $Y_1$ 和 $Y_2$ 作多元多重回归.

\noindent (i) 提出并拟合适当的线性回归模型.

\noindent (ii) 分析残差.

\noindent (iii) 对 (a)(iii) 中的预测变量值构造两个响应变量的 95\% 预测椭圆. 将此椭圆与 (a) 和 (b) 中的预测区间作比较, 并作出评论.

\begin{mvtable}
\phantomsection\label{tab:7.6}\mvtabcap{表 7.6\quad 抗忧郁剂数据}
\footnotesize
\mvfit{%
\begin{tabular}{ccccccc}
\toprule
$y_1$&$y_2$&$z_1$&$z_2$&$z_3$&$z_4$&$z_5$\\
TOT&AMI&GEN&AMT&PR&DIAP&QRS\\
\midrule
3\,389&3\,149&1&7\,500&220&0&140\\
1\,101&653&1&1\,975&200&0&100\\
1\,131&810&0&3\,600&205&60&111\\
596&448&1&675&160&60&120\\
896&844&1&750&185&70&83\\
1\,767&1\,450&1&2500&180&60&80\\
807&493&1&350&154&80&98\\
1\,111&941&0&1\,500&200&70&93\\
645&547&1&375&137&60&105\\
628&392&1&1050&167&60&74\\
1\,360&1\,283&1&3\,000&180&60&80\\
652&458&1&450&160&64&60\\
860&722&1&1\,750&135&90&79\\
500&384&0&2\,000&160&60&80\\
781&501&0&4\,500&180&0&100\\
1\,070&405&0&1\,500&170&90&120\\
1\,754&1\,520&1&3\,000&180&0&129\\
\bottomrule
\end{tabular}
}
\end{mvtable}

\noindent{\footnotesize 资料来源: 见文献[24].}

\medskip
\noindent 7.26\quad \hyperref[tab:7.7]{表 7.7} 是一组关于纸浆纤维及其所造纸张的特性度量数据(也可见文献[19]或网页: www.prenhall.com/statistics. 纸浆纤维特性变量为: $z_1=$ 纤维算术长度, $z_2=$ 长纤维比例, $z_3=$ 细纤维比例, $z_4=$ 零伸展抗拉应力. 纸的特性为: $y_1=$ 折断长度, $y_2=$ 弹性模量, $y_3=$ 失效应力, $y_4=$ 爆裂强度. 总计 $n=62$ 个观测.

\noindent (a) 对响应变量 $Y_1,Y_2,Y_3$ 和 $Y_4$ 分别作回归分析.

\noindent (i) 拟合并提出适当的回归模型.

\noindent (ii) 分析残差. 检查具有高杠杆率的离群点或观测.

\noindent (iii) 给定 $z_1=0.330$, $z_2=45.500$, $z_3=20.375$, $z_4=1.010$, 对变量 SF$(Y_3)$ 构造 95\% 预测区间.
